\subsection{A coefficient-support obstruction: Problem 14.22}
\label{cand:14.22}
\problemstatement{14.22}

Write the additive rational group multiplicatively as
\(G=\{a_q:q\in\Q\}\), with \(a_qa_r=a_{q+r}\).
It is torsion-free and linear via \(a_q\mapsto U(q)\)
in \(\SL_2(\Q)\). The solution set of
\(S=\{[x,a_1]=1\}\) is all of \(G\), and is irreducible
in the group-equation Zariski topology. Indeed, evaluating
a word \(w\in G*\langle x\rangle\) at \(x=a_t\)
gives \(a_{m(w)t+b(w)}\), where \(m(w)\in\Z\)
is its exponent sum and \(b(w)\in\Q\) its coefficient sum.
An equation therefore has zero, one, or all points as
solutions. Proper closed sets in this one-variable topology
are finite, and \(G\) is infinite.

Let \(T\) be any finite system equivalent to \(S\).
Every word in \(T\) has \(m(w)=b(w)=0\). Choose \(d>0\)
such that all its coefficients lie in the cyclic group
\(C=\langle a_{1/d}\rangle\). Form
\[
 K=G*_C(C\times\langle x\rangle)
\]
and the natural epimorphism
\(\pi:G*\langle x\rangle\to K\).
Each word of \(T\) vanishes in the second factor, since
both exponent sums are zero. The group \(K\) is torsion-free:
normal form in an amalgamated free product conjugates a
torsion element into a factor, and both factors are
torsion-free. Thus \(\ker\pi\) is a normal isolated
subgroup containing \(T\), and contains the least such
subgroup \(\sqrt T\).

But \([x,a_{1/(2d)}]\) vanishes at every point of \(G\),
so belongs to \(\operatorname{Rad}_G(T)\). Its image in
\(K\) is a reduced alternating word of length four:
\(x\notin C\) and \(a_{1/(2d)}\notin C\). It is nontrivial
by normal form. Hence
\[
 \sqrt T\ <\ \operatorname{Rad}_G(T)
\]
for every finite equivalent \(T\). The radical is isolated
because it is an intersection of kernels of maps into a
torsion-free group, which also justifies the displayed
inclusion.

The coefficient group is not finitely generated, as permitted
by the printed question. The separating quotient need only
be torsion-free; isolatedness does not mean unique roots.
In particular, commuting with a square does not imply
commuting with its root in that quotient.

For completeness, the archive also avoids the possible concern
that the intended geometry uses nonabelian coefficients.
Take \(G=\Q*\langle b\rangle\), represented faithfully by
\(a_q\mapsto U(q)\), \(b\mapsto
\left(\begin{smallmatrix}1&0\\t&1\end{smallmatrix}\right)\)
over \(\Q(t)\). For a cyclically reduced alternating word,
the trace has nonzero leading coefficient
\(\prod_j q_j n_j\), where the nonzero rational syllables
are \(q_j\) and the \(b\)-syllables are \(b^{n_j}\).
Thus the representation is faithful. The solution set of
\([x,a_1]=1\) is the factor \(\Q\), since free factors are
malnormal.

The coordinate group is
\((\Q\times\langle x\rangle)*\langle b\rangle\).
To verify this and irreducibility directly, represent \(x\)
by \(U(z)\) with \(z\) a second indeterminate. The same
leading-coefficient argument is faithful. Every nontrivial
word remains nontrivial under all but finitely many rational
specializations \(z=q\): a nonzero coefficient polynomial
has only finitely many roots. These evaluations discriminate
the coordinate group and make every proper relative algebraic
subset of the solution set finite.

For a finite equivalent system, choose a common denominator
for all rational syllables in its coefficients. Its words
already vanish in
\((C\times\langle x\rangle)*\langle b\rangle\), which
embeds in the coordinate group by normal form. They
therefore vanish in the torsion-free group
\[
 \bigl(\Q*_C(C\times\langle x\rangle)\bigr)*\langle b\rangle,
\]
while \([x,a_{1/(2d)}]\) again survives. This gives the same
strict radical obstruction with nonabelian linear coefficients.
Source: \artifact{research/14.22-proof.md}.
